\chapter{Current Results}
\label{ch:results}

\section{Summary of Progress}

\begin{table}[h]
\centering
\caption{Chronological progress from baseline to current best}
\label{tab:progress}
\begin{tabular}{clcc}
\toprule
\textbf{Step} & \textbf{Change} & \textbf{PCC} & \textbf{Gain} \\
\midrule
0 & L1 loss, fully-connected, from scratch & 0.467 & --- \\
1 & + Huber + Ranking loss & 0.486 & +0.019 \\
2 & + $k$-NN GAT ($k=30$) & \textbf{0.510} & +0.024 \\
\midrule
\multicolumn{2}{l}{\textbf{Total improvement}} & & \textbf{+0.043} \\
\bottomrule
\end{tabular}
\end{table}

\section{Comparison with Related Work}

\begin{table}[h]
\centering
\caption{Comparison with published methods on MegaScale/PNAS benchmark}
\label{tab:comparison}
\begin{tabular}{lcc}
\toprule
\textbf{Method} & \textbf{PCC (approx.)} & \textbf{Uses Structure?} \\
\midrule
BLOSUM62 baseline & $\sim$0.30 & No \\
ESM-2 zero-shot (log-likelihood ratio) & $\sim$0.45 & No \\
DeepEF (this work, current) & 0.510 & Yes (coordinates) \\
DeepEF (original, with pretraining) & $\sim$0.54 & Yes \\
RaSP & $\sim$0.60 & Yes \\
ThermoMPNN & $\sim$0.75 & Yes \\
\bottomrule
\end{tabular}
\end{table}

Note: Comparisons are approximate as methods use different test splits and evaluation protocols.

\section{What We Have vs.\ What We Need}

\begin{table}[h]
\centering
\caption{Current state assessment}
\label{tab:state}
\begin{tabular}{p{6cm}p{6cm}}
\toprule
\textbf{What We Have (Completed)} & \textbf{What We Need (Remaining)} \\
\midrule
\vspace{0.1cm}
$\checkmark$ Working training pipeline (from scratch) \newline
$\checkmark$ Huber + Ranking loss (validated) \newline
$\checkmark$ $k$-NN GAT architecture ($k=30$) \newline
$\checkmark$ Controlled ablation methodology \newline
$\checkmark$ PCC = 0.510 on 28-protein test set \newline
$\checkmark$ MegaScale data preprocessed \newline
$\checkmark$ AlphaFold PDB files identified (Zenodo) \newline
$\checkmark$ Code supports CLI flags for all experiments
&
\vspace{0.1cm}
$\square$ Training improvements (cosine LR, grad accum) \newline
$\square$ SaProt embeddings (1280-dim, structure-aware) \newline
$\square$ Download PDBs + install Foldseek \newline
$\square$ Per-protein fine-tuning on test proteins \newline
$\square$ Multi-seed ensemble (5 seeds) \newline
$\square$ Embedding projection into GNN branches \newline
$\square$ Final PCC target: $\geq$ 0.70
\\
\bottomrule
\end{tabular}
\end{table}

\section{Analysis of Current Limitations}

\subsection{Gap to Target}

Current PCC = 0.510, target $\geq$ 0.70. The gap of +0.19 must come from:
\begin{enumerate}
    \item \textbf{Input quality:} ProtT5 embeddings are sequence-only (Spearman $\sim$0.65). SaProt is structure-aware (Spearman 0.724). This single change should provide +0.03--0.05 PCC.
    \item \textbf{Training suboptimality:} Only 15 epochs with constant LR. Cosine schedule + more epochs + gradient accumulation = smoother convergence.
    \item \textbf{No pretraining:} The original DeepEF used native/decoy pretraining. Without it, the model must learn both structural representations and stability prediction simultaneously.
    \item \textbf{Single model variance:} A single random seed has high variance. Ensembling 5 seeds reduces noise by $\sim\sqrt{5}$.
\end{enumerate}

\subsection{What Cannot Be Fixed}

\begin{itemize}
    \item \textbf{No pretrained checkpoint:} The 43\_final\_model.pt from CASP12 pretraining is unavailable. This likely accounts for $\sim$0.08--0.12 PCC gap.
    \item \textbf{Dataset size:} 340 training proteins with $\sim$50K mutations is sufficient but not as large as ThermoMPNN's training set.
    \item \textbf{Hardware constraint:} 8GB VRAM limits batch size and model scale.
\end{itemize}
